% ch17 — Binary Black Hole Simulations
\chapter{Binary Black Hole Simulations}
\label{ch:binary-black-holes}

This chapter brings together the machinery of the preceding chapters ---
initial data, BSSN/CCZ4 evolution, AMR, gauge conditions, wave extraction,
and horizon finding --- to simulate the inspiral, merger, and ringdown of
binary black hole systems. It serves as both a capstone for Part~IV and a
practical guide to running and validating full numerical relativity
simulations.

\section{Orbital Dynamics and the PN--NR Interface}
\label{sec:bbh-orbits}

We review the three phases of binary evolution --- inspiral, merger,
ringdown --- and discuss where post-Newtonian (PN) approximations hand off
to full numerical evolution, motivating the need for NR simulations.

\section{Binary Puncture Initial Data}
\label{sec:bbh-initial-data}
\coderef{initial\_data}{rs}

We construct Bowen--York initial data for two boosted, spinning punctures,
compute the ADM mass and momenta, and discuss quasi-circular orbit
parameters from effective-potential methods.

\section{Evolution Setup and Gauge Choices}
\label{sec:bbh-evolution}
\coderef{bssn}{rs}
\coderef{gauge}{rs}

We configure the BSSN or CCZ4 evolution with 1+log slicing and
Gamma-driver shift for a binary system, discussing puncture tracking,
AMR box placement around each black hole, and Courant factor selection.

\section{Merger Dynamics}
\label{sec:bbh-merger}

We describe the phenomenology of the plunge and merger: common apparent
horizon formation, the transition from two individual horizons to one,
and the ringdown to a final Kerr black hole characterised by mass and
spin.

\section{Convergence Testing and Richardson Extrapolation}
\label{sec:bbh-convergence}

We demonstrate self-convergence of the waveform at multiple resolutions,
apply Richardson extrapolation to estimate the continuum limit, and
compute the convergence order to validate the numerical scheme.

\section{Waveform Analysis}
\label{sec:bbh-waveform}
\coderef{wave\_extraction}{rs}
\coderef{WaveExtraction}{hs}

We extract the dominant $(2,2)$ mode of $\Psi_4$, integrate to obtain
strain, compute radiated energy and final spin, and compare with
semi-analytic predictions from the test-particle limit and surrogate
models.

\section{Parameter Space: Mass Ratio and Spin}
\label{sec:bbh-parameter-space}

We survey how the dynamics and waveforms change across the binary
parameter space --- mass ratio $q$, spin magnitudes $\chi_i$, and spin
orientations --- highlighting precession, recoil kicks, and spin-orbit
coupling.

\paragraph{Key References.}
The moving-puncture breakthrough follows \cite{Campanelli2006}; initial
data and quasi-circular parameters draw on \cite{Cook2000,Bowen1980};
convergence testing follows \cite{Baumgarte2010}; the CCZ4 formulation
follows \cite{Alic2012}.
